% !TEX program = xelatex
\documentclass[10pt,letterpaper,twocolumn]{article}

\usepackage[top=0.75in, bottom=0.75in, left=0.75in, right=0.75in]{geometry}
\usepackage{fontspec}
\usepackage{unicode-math}
\setmainfont{Latin Modern Roman}[
  Extension=.otf,
  UprightFont=lmroman10-regular,
  BoldFont=lmroman10-bold,
  ItalicFont=lmroman10-italic,
  BoldItalicFont=lmroman10-bolditalic,
]
\setsansfont{Latin Modern Sans}[
  Extension=.otf,
  UprightFont=lmsans10-regular,
  BoldFont=lmsans10-bold,
  ItalicFont=lmsans10-oblique,
  BoldItalicFont=lmsans10-boldoblique,
]
\newfontfamily\acbold{ywft-processing-bold}[Path=../../system/public/type/webfonts/,Extension=.ttf]
\newfontfamily\aclight{ywft-processing-light}[Path=../../system/public/type/webfonts/,Extension=.ttf]
\setmonofont{Latin Modern Mono}[
  Extension=.otf,
  UprightFont=lmmono10-regular,
  ItalicFont=lmmono10-italic,
  Scale=0.85,
]

\usepackage{xcolor}
\usepackage{titlesec}
\usepackage{enumitem}
\usepackage{booktabs}
\usepackage{tabularx}
\usepackage{fancyhdr}
\usepackage{hyperref}
\usepackage{xurl}
\usepackage{graphicx}
\graphicspath{{figures/}{../../papers/arxiv-ac/figures/}}
\usepackage{ragged2e}
\usepackage{microtype}
\usepackage{natbib}
\usepackage[colorspec=0.92]{draftwatermark}

\definecolor{acpink}{RGB}{180,72,135}
\definecolor{acpurple}{RGB}{120,80,180}
\definecolor{acdark}{RGB}{64,56,74}
\definecolor{acgray}{RGB}{119,119,119}
\definecolor{draftcolor}{RGB}{180,72,135}

\DraftwatermarkOptions{text=WORKING DRAFT,fontsize=3cm,color=draftcolor!18,angle=45}

\hypersetup{colorlinks=true,linkcolor=acpurple,urlcolor=acpurple,citecolor=acpurple,
  pdftitle={Sucking on the Complex: Platformhegemoni, kritik-som-indhold og behovet for anti-miljøer}}

\titleformat{\section}{\normalfont\bfseries\normalsize\uppercase}{\thesection.}{0.5em}{}
\titlespacing{\section}{0pt}{1.2em}{0.3em}
\titleformat{\subsection}{\normalfont\bfseries\small}{\thesubsection}{0.5em}{}
\titlespacing{\subsection}{0pt}{0.8em}{0.2em}

\pagestyle{fancy}\fancyhf{}
\renewcommand{\headrulewidth}{0pt}
\fancyhead[C]{\footnotesize\color{acpink}\textit{Arbejdsudkast --- ikke til citation}}
\fancyfoot[C]{\footnotesize\thepage}

\newcommand{\ac}{\textsc{Aesthetic.Computer}}
\newcommand{\np}{\textsc{notepat}}

\setlist[itemize]{nosep, leftmargin=1.2em, itemsep=0.1em}
\setlength{\columnsep}{1.8em}
\setlength{\parindent}{1em}
\setlength{\parskip}{0.3em}

\begin{document}

\twocolumn[{%
\begin{center}

\includegraphics[height=4em]{pals}\par\vspace{0.5em}
{\acbold\fontsize{22pt}{26pt}\selectfont\color{acdark} Sucking on the Complex}\par
\includegraphics[width=0.5\textwidth]{figures/cover}\par
\vspace{0.35em}
\vspace{0.2em}
{\aclight\fontsize{11pt}{13pt}\selectfont\color{acpink} Platformhegemoni, kritik-som-indhold og behovet for anti-milj{\o}er}\par
\vspace{0.6em}
{\normalsize\href{https://prompt.ac/@jeffrey}{@jeffrey}}\par
{\small\color{acgray} Aesthetic.Computer}\par
{\small\color{acgray} ORCID: \href{https://orcid.org/0009-0007-4460-4913}{0009-0007-4460-4913}}\par
\vspace{0.3em}
{\small\color{acpurple} \url{https://aesthetic.computer}}\par
\vspace{0.6em}
\rule{\textwidth}{1.5pt}
\vspace{0.5em}
\end{center}

\begin{center}
{\small\color{acpink}\textbf{[ arbejdsudkast --- ikke til citation ]}}
\end{center}
\vspace{0.3em}

\begin{quote}
\small\noindent\textbf{Abstrakt.}
Den samtidige kunstverdens medieinfrastruktur domineres af to virksomheder: Meta og ByteDance. For de fleste arbejdende kunstnere er en platformtilstedeværelse en forudsætning for kulturel læselighed, og kritikken af platforme er blevet et af de mest fejrede indholdsformater, platformene selv er vært for. Et tidligere udkast af denne artikel argumenterede for, at \emph{anti-miljøer}---McLuhans begreb for modstrukturer bygget uden for et dominerende miljø---er den eneste flugt fra denne fælde, og gennemgik fremtrædende kunstnere, hvis platformkritiske værker absorberes af de institutioner, der udstiller dem. Det argument er stadig her, komprimeret. Det nuværende udkast inverterer artiklens tyngdepunkt: i stedet for at gennemgå diagnosen af platformhegemoni eller den strukturelle fiasko for institutionel kritik-som-indhold, rapporterer jeg fra fem års arbejde med at bygge et alternativ---\ac{}~\citep{scudder2026ac}, en runtime, et socialt netværk, et instrument og et operativsystem. Det, som dette arbejde understøtter, er en smallere og mere konkret påstand end total afvisning. Platforme er hverken modstander eller hjem; de er distributionsmaskiner. Det interessante træk for en kunstner i 2026 er ikke at afvise maskinen, men at bruge den bevidst---og levere en nyttelast af grafisk notation og interface-retorik, som platformen ikke kan inspicere, mod en \emph{brugerdom}, som platformen ikke kan rumme. Whistlegraph~\citep{scudder2026whistlegraph} er sagen. \ac{} er det nedstrøms. Jeg katalogiserer seks strukturelle forpligtelser, som et anti-miljø for kreativ databehandling hviler på---URL-først-adressering, stykke-ikke-opslag som afgrænset enhed, instrument-ikke-profil som grænseflade, et kompositions-æra-sprog, selvstændige noder som hardware, og affinitet-ikke-følgeskab som fællesskab---og argumenterer for, at arbejdet med at bygge et anti-miljø ikke er arbejdet med at afvise platformen, men at prime en brugerdom, som platformen kan distribuere til, men ikke rumme.
\end{quote}
\vspace{0.5em}
}]

\section{Maskinen og det ydre}

Standardrammen for platformkritik siger: komplekset er totaliseret, kunstneren, der fodrer det, er medskyldig, og det ærlige træk er at afvise det. Denne artikel blev oprindeligt skrevet i det register. Den nuværende version er ikke. Fem års arbejde med at bygge \ac{}~\citep{scudder2026ac} sammen med et årtis drift af @whistlegraph~\citep{scudder2026whistlegraph} har lært, at registret er forkert på en bestemt måde: det fejllæser, hvad platforme er \emph{til} for, set fra kunstnerens side.

Platforme er maskiner. De flytter opmærksomhed med en hastighed og i et volumen, som ingen institutionel pipeline kommer i nærheden af, og de gør det pålideligt nok til, at en kunstner kan behandle dem som infrastruktur på samme måde som man behandler elektricitet eller en postrute. En kunstner, der bekymrer sig om at bygge en praksis, som overlever platformen, har et valg, som kritiklitteraturen næsten aldrig beskriver: ride maskinen bevidst og bruge dens rækkevidde til at prime et publikum i en formel grammatik, som platformen hverken har forfattet eller kan reproducere, og rute dette publikum mod en praksis og en \emph{brugerdom}, der lever på en anden stak. Dette er ikke afvisning. Det er tættere på ekstraktion: at bruge røret til at levere en nyttelast, som røret ikke kan inspicere.

Whistlegraph var en sådan nyttelast. Formen opstod af Radical Digital Painting~\citep{scudder2017manifesto}, 65 præ-pandemiske forelæsnings-performances inklusive Goodiepal \& Pals' europæiske turné i 2018~\citep{scudder2026goodiepal} og 35. Chaos Communication Congress~\citep{scudder2018rdp35c3}, og et årtis tegne-som-sang-praksis, der gik forud for kontoen. TikTok gjorde ikke whistlegraphen viral; whistlegraphen gjorde sig selv viral på TikTok, fordi formens strukturelle egenskab---reproducérbarhed, mark-for-stavelse-korrespondance, partituret der lærer sig selv~\citep{scudder2026whistlegraph}---allerede var der. Kontoens 2,6 millioner følgere er ikke værket. De er det primede publikum for det, der kommer efter formens distributionsfase: en \emph{brugerdom}---en befolkning organiseret omkring et kreativt databehandlingssystem, de selv kan køre, forke og optræde inde i, snarere end et publikum organiseret omkring indhold, som de kun kan scrolle.

Det følgende er hverken en gennemgang af platformhegemoni (den diagnose er velkendt; \S2 komprimerer den) eller en frisk anklage mod institutionel kritik-som-indhold (\S3). Det er en redegørelse for, hvad fem års bygning af den nedstrøms stak har produceret, hvilke strukturelle forpligtelser et anti-miljø for kreativ databehandling hviler på, hvordan maskinen primer en brugerdom til det, og hvad det koster at holde cirklen åben. Artiklen, du læser, er selv en af disse forpligtelser. Den er skrevet indefra cirklen, citerer resten af cirklens dokumentation~\citep{scudder2026ac, scudder2026kidlisp, scudder2026notepat, scudder2026os, scudder2026pieces, scudder2026urltradition, scudder2026archaeology, scudder2026whistlegraph, scudder2026goodiepal}, og distribueres gennem den forskningsplatter, som cirklen har frembragt.

\section{Komplekset}

Meta (Facebook, Instagram, Threads, WhatsApp) og ByteDance (TikTok, Douyin) medierer tilsammen den kulturelle synlighed for næsten enhver arbejdende kunstner på jorden. Instagram alene har over tre milliarder månedligt aktive brugere; det er den primære opdagelsesplatform for gallerier, kuratorer, samlere og institutioner. TikToks kortformat-video er blevet standarddistributionsmekanismen for performance, atelierpraksis og kunstuddannelse. Tilsammen udgør de, hvad \citet{nieborg2018platformization} kalder ``platformiseringen af kulturel produktion'': den kontingente kulturelle vare, formet ikke af kunstnerisk intention, men af platformens sorteringslogik. For yngre og midtkarriere-kunstnere er platformmålinger---følgere, engagementsrate, postfrekvens---blevet stedfortrædende mål for kulturel relevans. Prisen for synlighed er data: enhver interaktion fodrer reklamemaskinen, der finansierer platformen~\citep{zuboff2019surveillance, srnicek2017platform}.

Algoritmen viser ikke passivt kunst. Den former aktivt, hvilken kunst der bliver lavet. Når en kunstner poster et maleri, og det modtager højt engagement, viser algoritmen det til flere mennesker; kunstneren, der ser responsen, laver flere malerier som det. Dette er ikke censur---det er adfærdsmæssig påvirkning, den centrale mekanisme, som \citet{zuboff2019surveillance} identificerer som overvågningskapitalismens kerne: forudsigelse og modifikation af menneskelig adfærd i stor skala, for profit. Og på Instagram optager et Richter-maleri den samme 1080-pixel firkant som et fotografi af nogens frokost. Platformen skelner ikke mellem årtiers arbejde og sekunders arbejde. Dette er ikke en fejl i kurateringen; det er platformens design~\citep{chun2016updating}. Engagement er den eneste måling. \citet{lovink2011networks} kalder den resulterende struktur ``netværket uden en sag'': et system, der forbinder alt og kontekstualiserer intet. \citet{noble2018algorithms} og \citet{benjamin2019race} viser, at det rangerende apparat under udfladningen ikke er neutralt; det indkoder og forstærker eksisterende hierarkier. \citet{terranova2000free} identificerede konsekvensen for femogtyve år siden som ``frit arbejde'': ulønnet kulturel produktion, der genererer værdi for platforme. Kunstnere er opmærksomhedsøkonomiens mest dedikerede friarbejdere.

Den dybeste ramme for dette er Attalis. I \emph{Noise}~\citep{attali1985noise} argumenterer han for, at koncertsalens orkester er et instrument for statsmagt: det arrangerer kroppe i hierarkiske positioner under en dirigents autoritet, kræver synkron lydighed mod et skrevet partitur og producerer et samlet output fra kontrolleret individuelt arbejde. Dette er musikkens \emph{repræsentations}-æra---musik som spektakel, som bevis på at magt kan organisere støj til harmoni. Teknologiplatformen er det enogtyvende århundredes orkester. Grundlæggeren er dirigenten. Algoritmen er partituret. Brugere, skabere og kunstnere er musikerne---der hver spiller deres rolle, hver tror de udtrykker sig selv, og hver producerer værdi, der tilfalder institutionen. Platformens retorik om at ``give alle en stemme'' er repræsentations-æra-sprog klædt i kompositions-æra-tøj. Alle har en stemme, men algoritmen bestemmer, hvem der bliver hørt. Attali glorificerede det, der skulle komme efter repræsentationen---\emph{kompositionen}, en tilstand hvor skellet mellem udøver og publikum opløses, hvor musik laves for glæden ved at lave den. Platformen er ikke komposition. Den er repræsentation i planetarisk skala, med den yderligere udvinding af adfærdsdata, som selv hoforkestret aldrig opnåede. Tech-startuppet er strukturelt tættere på Habsburgernes hoforkester end på folkesangen. Forskellen er, at hoffet betalte sine musikere.

Diagnosen er ikke original og behøver ikke at være det. Det følgende er.

\section{Kritik-som-indhold}

Det dominerende kunstneriske svar på platformhegemonien har været at lave værker \emph{om} den. Disse værker udstilles i verdens mest prestigefyldte kulturinstitutioner. De fejres, samles og diskuteres. De strukturelle forhold, de beskriver, forbliver uforandrede.

Modaliteten har flere former. Paglen fotograferer klassificerede militære installationer og dokumenterer overvågningens fysiske infrastruktur, udstillet på Whitney, MoMA, Smithsonian og Venedig-biennalen; værket reklamerer for kompleksets rækkevidde uden at rute sit publikum mod noget andet. \citet{steyerl2017duty} beskriver en verden af ``toldfri kunst'' der cirkulerer i frihandelszoner, og hendes videoinstallationer på Serpentine, Venedig-biennalen og Park Avenue Armory er analyser af den samme cirkulation---delt på Instagram, diskuteret på TikTok, kritikken distribueret af den maskine, den beskriver (\citet{steyerl2009poor} argumenterede for, at det ``fattige billede'' får politisk kraft gennem sin lavopløsnings-cirkulation; i feedet genererer hvert billede den samme adfærdsdata). \citet{blas2014informatic} udviklede ``Facial Weaponization Suite'' (2011--2014), kollektive masker, der besejrer ansigtsgenkendelse, vist på Whitechapel og Tate Modern---men masken er et kunstobjekt i en glasmontre, ikke et værktøj, man kan bære. Mindy Seus \emph{Cyberfeminism Index}~\citep{seu2023cyberfeminism} katalogiserer over 700 ressourcer, der sporer cyberfeministisk modstand, meget af det distribueret gennem Instagram Stories: et feministisk arkiv over digital modstand, gjort læseligt for kunstverdenen gennem Metas overvågningsinfrastruktur.

\citet{ulman2014excellences} er det kanoniske tilfælde. Hendes fem måneder lange Instagram-performance i 2014---en skriptet manipulation af platformens logik om stræben og selvfremstilling, senere udstillet på Tate Modern i 2016---lykkedes netop fordi den var umulig at skelne fra indhold. Platformen vidste ikke, at den blev brugt som materiale. Adfærdsdataen, den genererede, var identisk med den fra ægte livsstilsopslag. Dette er den strukturelle pointe, ikke en anklage mod nogen individuel kunstner: algoritmen skelner ikke mellem kritik og billigelse. Paglens fotografi af en NSA-facilitet og et landskabsmaleri genererer begge engagement; begge fodrer maskinen. Pipelinen er: lav indhold til feedet $\rightarrow$ bliv opdaget af kuratorer (der fandt dig på Instagram) $\rightarrow$ udstil det samme indhold eller en kritik af det i institutioner, hvis publikum også blev fundet på Instagram. \citet{fraser2005critique} identificerede denne sløjfe i institutionel kritik for tre årtier siden; platformhegemoni er det samme strukturelle problem med overvågningskapitalisme lagt oven på.

Den specifikke fiasko ved kritik-som-indhold er ikke, at den opererer inden for platformen---det gør alt---men at den \emph{kritiserer uden at dyrke}. Den bruger platformens rækkevidde til at forstærke en diagnose og stopper der. Den bruger ikke rækkevidden til at prime et publikum for en praksis, der lever på en anden stak. Platformen er fuldstændig villig til at distribuere diagnosen; diagnosen er indhold, og indhold er, hvad platformen distribuerer. Det, platformen ikke kan distribuere, er en brugerdom. Det er arbejde, der skal gøres bevidst, med en form, som platformen kan bære, men ikke reproducere.

\section{Frie rum for leg}

Kunst vokser ikke i feeds. \citet{winnicott1971playing} kaldte det rum, hvor kreativ erfaring bliver mulig, \emph{potentielt rum}---et overgangsområde mellem indre og ydre virkelighed, skrøbeligt, der kræver sikkerhed, kontinuitet og frihed fra indtrængen. \citet{huizinga1938homo} beskrev legens ``magiske cirkel'': et afgrænset rum med sine egne regler, adskilt fra hverdagslivet. Inde i den magiske cirkel er ydreverdenens regler suspenderet. Et atelier, når det fungerer, er en magisk cirkel. Et øveværelse er en magisk cirkel. En skitsebog er en magisk cirkel.

Platformen har ingen magisk cirkel. Enhver kreativ handling udført på Instagram bliver øjeblikkeligt målt, sorteret, rangeret og fodret tilbage som engagementsdata. Det potentielle rum kollapser til målinger, i det øjeblik værket er postet. \citet{fisher2009capitalist} kaldte den gennemgribende fornemmelse af, at der ikke findes noget alternativ, \emph{kapitalistisk realisme}: realismen er ikke i indholdet, men i manglende evne til at forestille sig andet. Kunstskolen underviser i den magiske cirkels form; platforme gør det ikke. Denne artikel er et argument for at forestille sig andet---og en rapport om, hvordan ``andet'' ser ud, når det er bygget.

\section{Anti-milj{\o}er}

McLuhan argumenterede for, at man ikke kan opfatte et miljø indefra~\citep{mcluhan1964understanding, mcluhan1969playboy}. Et miljø er usynligt for dets beboere, netop fordi det er totalt. Den eneste måde at gøre et miljø synligt på er at konstruere et \emph{anti-miljø}: en modstruktur, hvis forskel fra det dominerende miljø afslører miljøets form---hvad \citet{shklovsky1917art} kaldte \emph{defamiliarisering}, kunstens evne til at gøre det vante fremmed.

Denne artikel bruger \emph{anti-miljø} i dets strukturelle betydning: ikke en protest, ikke et udstillingsvindue, ikke en afvisning. Et arbejdende alternativ. Et instrument. En afgrænset, regelbærende struktur, der er bygget sideløbende med det dominerende miljø og router kreativt arbejde gennem sin egen stak snarere end gennem platformens. Forholdet til platformen er hverken oppositionelt eller afholdende. Et anti-miljø kan bruge en platforms rækkevidde uden at arve dens logik, på samme måde som mail-kunst brugte postvæsenet uden at blive et postvæsen. Lokalradio, mail-kunst-netværk, zine-kultur, Fluxus' distributionskredsløb og Nelsons \emph{Computer Lib}~\citep{nelson1974computerlib} fungerede alle som anti-miljøer for deres tids broadcast-medier; \citet{lanier2018ten} når en beslægtet konklusion indefra industrien. Nogle overlevede; mange blev absorberet. Overlevelsen afhang mindre af afvisning end af, om anti-miljøet havde en praksis, som det dominerende miljø ikke kunne reproducere.

Et anti-miljø for kreativ databehandling i 2026 er en stak: en runtime, et sprog, et instrument, et hardware-substrat, en fællesskabsprotokol. Hvert lag er en specifik strukturel forpligtelse, der afviser platformlogik på laggets niveau, ikke på budskabets. Næste afsnit beskriver de forpligtelser, \ac{} har indgået, hver enkelt dokumenteret i dybden af en ledsagende artikel i den medfølgende forskningsplatter.

\section{Hvad der er blevet bygget}

Fem års drift har produceret et sæt strukturelle forpligtelser. Hver er en specifik afvisning af platformlogik på stakkens niveau, og hver er dokumenteret i dybden af en søskende-artikel i den medfølgende forskningsplatter. Dette afsnit er syntesen: disse er ikke adskilte beslutninger. De er det samme træk.

\subsection{URL-f{\o}rst: adressen er gr{\ae}nsefladen}

Hvert stykke på \ac{} er en URL. Der er intet feed, ingen menu, ingen filvælger, ingen projektliste, ingen algoritmisk overflade. Prompten \emph{er} adresselinjen: man skriver et ord, og man er i stykket~\citep{scudder2026urltradition}. Dette er HyperCard-traditionen genfundet i nettet, nedstigende gennem Glitchs øjeblikkelige remix-links og p5.js' Web Editor-URL'er. En platform har en hjemmeside, man ankommer til; \ac{} har hundrede adresser, man skriver mod. URL'en er ikke en funktion. Den er hele systemets mediumegenskab.

\subsection{Stykke-ikke-opslag: den afgr{\ae}nsede enhed}

Enheden for kulturel produktion på \ac{} er et \emph{stykke}: et selvstændigt interaktivt program~\citep{scudder2026pieces}, adresserbart ved navn, vedvarende over tid, ikke algoritmisk sorteret. Et stykke er for et opslag, hvad en sang er for et scroll. Man opdager ikke et stykke, fordi en algoritme overfladede det; man ankommer til et stykke, fordi nogen navngav det for en, eller fordi man skrev dets adresse, eller fordi man fandt det i den offentligt søgbare liste. Winnicotts \emph{potentielle rum}~\citep{winnicott1971playing} er genoprettet på enhedens niveau: hvert stykke er afgrænset, regelsat og uforstyrret af engagementsmålinger.

\subsection{Instrument-ikke-profil: gr{\ae}nsefladen spilles}

\np{}~\citep{scudder2026notepat} er en kromatisk tastatur-synthesizer, der kører i browseren og på bare metal. Det er et instrument, man spiller. Det er ikke en profil, man kuraterer, ikke et feed, man scroller, ikke et brand, man vedligeholder. Forskellen mellem instrument og profil er forskellen mellem brug og fremvisning. En profil er optimeret til at blive set; et instrument er optimeret til at blive brugt. \ac{} behandler den kreative grænseflade som et instrument på alle niveauer af stakken---fra \np{}s tastatur til KidLisps REPL til selve prompten---og dette valg nedprioriterer læselighed-for-andre som designmål til fordel for umiddelbarhed-for-selvet.

\subsection{KidLisp: et kompositions-{\ae}ra-sprog}

KidLisp~\citep{scudder2026kidlisp} er en minimal Lisp-dialekt for generativ kunst. Det er et kompositions-æra-sprog i Attalis forstand~\citep{attali1985noise}: skellet mellem udøver og publikum opløses inde i det. Den samme kode, man skriver for at male et billede, er den kode, andre læser for at forstå billedet. Den samme REPL, man bruger til at skitsere, er den REPL, andre bruger til at forke. Dette er ikke en generel egenskab ved Lisp---de fleste Lisps er arkitekt-æra, bygget til at specificere systemer---men en egenskab ved KidLisp specifikt, som bytter ekspressiv kraft for umiddelbarhed, forkbarhed og læselighed ved 8-bit opløsning.

\subsection{Selvst{\ae}ndig node: hardwaren er din}

\ac{} leveres som et bare-metal-operativsystem, der booter overskuds-hardware direkte ind i den kreative runtime~\citep{scudder2026os, scudder2026identity}. Der er ingen cloud-konto, intet vendor-lock, ingen udløbstimer for opdateringer, ingen adfærdstelemetri. Den bærbare er en selvstændig node på netværket: den kører sin egen server, holder sine egne data, fødererer med andre noder på egne vilkår. Dette er anti-Chromebook'en~\citep{scudder2026openschools}: den samme 50-dollar-hardware, flashet med software, brugeren ejer, på en protokolstak, de kan granske. Med omkring 240 millioner Windows 10-maskiner ved deres udløb på vej til lossepladsen i 2025--2026-cyklussen, er den selvstændige-node-forpligtelse ikke en niche-politisk gestus; det er et materielt program for, hvad man skal gøre med hardwaren.

\subsection{Pals: affinitet, ikke f{\o}lgeskab}

\ac{}-fællesskabet er organiseret efter Pals-modellen, arvet fra Goodiepals turnéer og dokumenteret i dybden af Goodiepal-artiklen~\citep{scudder2026goodiepal}. En Pal er ikke en følger. En Pal er en, der laver med en, rejser med en, deler instrumentet med en. Forholdet er gensidigt, affinitetsbaseret og ansigt-til-ansigt, hvor det er muligt; det er ikke en envejs broadcast-kanal. Platforme optimerer for parasocial en-til-mange-strøm, fordi det er, hvad reklame monetariserer. Pals optimerer for en-til-en-dybde, fordi det er, hvad praksis belønner. Tallene er mindre. Båndene er tættere. Modellen skalerer ved at forke, ikke ved rækkevidde. Lewis' ``making kin with the machines''~\citep{lewis2018making}, \citet{escobar2018designs}s designs for pluriverset og \citet{costanzachock2020design}s design-retfærdighed artikulerer den samme forpligtelse i forskellige registre: relationelt snarere end transaktionelt, omsorg snarere end udvinding.

\subsection{94-projekt-linjen}

Intet af ovenstående er nyt. \ac{}-repositoriearkæologien sporer platformens tekniske genealogi gennem 94 forgængerprojekter fra 2007--2020~\citep{scudder2026archaeology}---tegneværktøjer, performance-værktøjer, udviklingsinfrastruktur, spil, instrumenter. Forpligtelserne ovenfor ankom ikke som en designfilosofi. De akkumulerede som løsningen på problemer, der blev ved med at vende tilbage gennem 94 forsøg. Anti-miljøet er ikke en idé, forfatteren havde. Det er den form, der faldt ud af at praktisere den samme praksis i to årtier inden for medie-miljøer, der hver i deres tur var totale.

\section{Maskinen og brugerdommen}

Indvendingen mod argumentet hidtil skriver sig selv: hvis anti-miljøer er bygget uden for platformen, hvad gør man så af @whistlegraph-TikTok-kontoen, der nåede 2,6 millioner følgere under COVID-19-pandemien~\citep{scudder2026whistlegraph} og stadig flytter omkring otteogfyrre tusind videovisninger om dagen mere end to år efter trioen gik i dvale i november 2023?\footnote{Øjebliksbillede fra TikTok Analytics for @whistlegraph, 24. april 2025 til og med 23. april 2026---det andet fulde år efter trioens opløsning, hvor ingen nye whistlegraphs blev produceret af de oprindelige skabere: $\sim$17,69 millioner videovisninger (median 43\,047/dag), $\sim$465\,000 likes, $\sim$16\,000 delinger, $\sim$191\,000 profilvisninger. Maksimale enkeltdagsvisninger i dette vindue (208\,000 den 27. september 2025; 198\,000 den 16. februar 2026) afspejler fortsat organisk cirkulation af det oprindelige arkiv.} Er det ikke bevis for, at det ydre kun eksisterer med platformens tilladelse?

Det er bevis for noget mere præcist. Platforme er distributionsmaskiner. De flytter opmærksomhed med en hastighed og i en skala, som ingen institutionel pipeline kommer i nærheden af, og pålideligt nok til, at en kunstner kan behandle dem som infrastruktur. Spørgsmålet, som standard-platformkritiklitteraturen ikke stiller---fordi den er forpligtet til en oppositionel ramme---er, hvad maskinen bruges til at \emph{bære}. @whistlegraph-kontoen bar en specifik formel grammatik: grafisk notation (mærker der svarer til sungne stavelser), interface-retorik (performance som forelæsning som instrument) og en reproducérbarhedslogik (partituret der lærer sig selv~\citep{scudder2026whistlegraph}). Den grammatik blev udviklet præ-platform gennem Radical Digital Painting~\citep{scudder2017manifesto, scudder2018rdp35c3}, 65 forelæsnings-performances inklusive en europæisk turné i 2018 med Goodiepal~\citep{scudder2026goodiepal}, og et årtis tegne-som-sang-praksis, som platformen ikke havde nogen del i at forfatte. TikTok skabte ikke whistlegraphen; den distribuerede én.

Det, distributionen primede, er en \emph{brugerdom}: en befolkning på omkring 2,6 millioner mennesker, der gennem at se og reproducere whistlegraphs lærte at genkende en specifik slags formelt træk---mærker som instruktioner, grænseflader som instrumenter, performance som pædagogik. Brugerdommen er ikke følgertællingen. Den er det primede publikum tilgængeligt nedstrøms, i den takt det kan konverteres til praksis på et system, der bærer den samme grammatik dybere, end platformen kan rumme. \ac{} er det system~\citep{scudder2026ac}. Brugerdommen er dets kimbefolkning. Netværksauditen~\citep{scudder2026audit} dokumenterer den nuværende form af den konverterede brugerdom---2.800+ registrerede håndtag, 16.000+ KidLisp-programmer, 93.000+ boot-sessioner---en lille brøkdel af det primede publikum og mere end noget sammenligneligt kreativt databehandlingsværktøj har produceret uden en institutionel forælder.

Paradokset, som det tidligere udkast af artiklen ikke kunne løse, er dette: det største anti-miljø i \ac{}-platteren blev bygget \emph{gennem} komplekset. Løsningen er, at komplekset aldrig var modstanderen. Det var maskinen. Fejlen, kritik-som-indhold-litteraturen begår, er at læse enhver brug af platformen som kapitulation. Kunstnerens faktiske træk---det, der er tilgængeligt for enhver med en form, som platformens algoritme vil forstærke---er at bruge maskinen bevidst og levere en nyttelast, platformen ikke kan inspicere, mod en brugerdom, platformen ikke kan rumme. Skalaen lejes af maskinen. Praksissen ejes af brugerdommen.

\section{Finansiering af maskinen og cirklen}

Anti-miljøer er svære at finansiere af en strukturel grund: platformøkonomien priser opmærksomhed, og anti-miljøet er bygget til at afvise opmærksomhedens målbarhed. Et stykke har ingen engagementsrate. En Pal har ingen cost-per-mille. Et instrument har ingen retention-kurve. Bæredygtighedsartiklen i denne serie~\citep{scudder2026sustainability} gennemgår, hvordan 28 sammenlignelige kreative databehandlingsværktøjer er blevet finansieret---akademisk subsidie, firmamæcenskab, personligt offer, crowdfunding i fællesskabet, uventet tilskud eller slet ingen model---og finder, at de værktøjer, der er mest brugt, ikke er de værktøjer, der er bedst finansieret. Den mediane afstand mellem et værktøjs skabelse og dets første bæredygtige finansiering er otte år. De fleste værktøjer lukker, før de krydser det.

Maskine-og-brugerdom-rammen løser ikke finansieringsproblemet, men den klargør det. Maskinen er gratis at bruge og gratis at forlade. Brugerdommen er det aktiv, maskinen producerer. Praksissen er, hvor aktivet bruges. Finansieringsspørgsmålet er ikke, om maskinen skal monetariseres---maskinen monetariserer kunstnere på sine egne betingelser og sælger deres opmærksomhed til annoncører---men hvordan man konverterer brugerdommen til en selvbærende praksis-økonomi, før maskinens næste algoritmeopdatering roterer opmærksomheden andetsteds hen. Illichs \emph{værktøjer for konvivialitet}~\citep{illich1973tools} navngiver bestemmelsesstedet: værktøjer bygget i menneskelig skala, med hånden, til brug snarere end forbrug. Økonomien for sådanne værktøjer er stadig under udarbejdelse. Denne artikel er én post i dens dokumentation.

\section{Koda}

Anti-miljøer er mulige. De er ikke sjældne; de er simpelthen ikke læselige for platformen, fordi platformen ikke blev bygget til at læse det, de producerer. Denne artikel har beskrevet formen for ét, indefra, med de specifikke strukturelle forpligtelser, det hviler på, skrevet ned, så andre kan genkende dem, forke dem, tilpasse dem eller afvise dem:

\begin{itemize}
  \item URL-først, ikke feed-først.
  \item Stykke-ikke-opslag: afgrænsede enheder, ikke scrollende strømme.
  \item Instrument-ikke-profil: grænseflader at spille, ikke at kuratere.
  \item Kompositions-æra-sprog: kode læst, som den er skrevet.
  \item Selvstændig node: hardware og data, du ejer.
  \item Affinitet, ikke følgeskab: skalering ved forking, ikke ved rækkevidde.
  \item Maskine, ikke modstander: brug platformens rækkevidde til at prime en brugerdom, hvis praksis lever andetsteds.
\end{itemize}

Dette er reglerne for én cirkel. Andre mennesker kan tegne deres egne. Algoritmen kan ikke læse nogen af dem.

\vspace{0.5em}
\noindent\textbf{ORCID:} \href{https://orcid.org/0009-0007-4460-4913}{0009-0007-4460-4913}

\vspace{0.5em}
\noindent\textit{Oversat fra engelsk. Originalversion tilg{\ae}ngelig p{\aa} \url{https://papers.aesthetic.computer}}

\bibliographystyle{plainnat}
\bibliography{references}

\end{document}
